\documentclass[11pt]{book}
\usepackage[paperwidth=145mm,paperheight=200mm,top=18mm,bottom=20mm,inner=20mm,outer=16mm]{geometry}
\usepackage{brumfiel}

\begin{document}
\setcounter{chapter}{2}

% ---- local macros (hoist into book preamble) ----
% Halftone photographic plate: framed placeholder, never redrawn (B&W pass).
% Carries the real plate since 2026-08-17 (Caleb: "I WANT THE ACTUAL HALFTONE
% PICTURES IN THERE"), cut from scan02 p.16 -- see plates/MANIFEST.md.
%   #1 = image file in plates/, #2 = name, #3 = caption
\newcommand{\IIplate}[3]{%
  \par\medskip
  \begin{center}
  \includegraphics[height=62mm]{#1}%
  \par\smallskip
  {\footnotesize\textsc{#2}\\ #3}
  \end{center}
  \par\medskip}
% Unnumbered star footnote. \starnote in brumfiel.sty emits a bare
% \footnotetext, which prints the footnote counter ("0"); the book marks this
% note with the same "*" the starred exercises carry, and no number.
\newcommand{\IIstarnote}{\begingroup\renewcommand{\thefootnote}{}%
  \footnotetext{*Starred exercises, sections, theorems, etc., are optional.}%
  \endgroup}
% Ch 2's eight "Problem." blocks. These are NOT the book's "Remark." blocks:
% p.22 etc. print an italic "Problem." head and nothing else, so wrapping them
% in {remark} (which emits "Remark.") printed a word the book never sets.
% Same shape as brumfiel.sty's {remark}, with the correct head.
\newenvironment{IIproblem}{\par\medskip\noindent\emph{Problem.}\ }{\par\medskip}
% ---- end local macros ----

% Facing-page plate for the chapter opener (book p.20, facing p.21).
\IIplate{plate-aristotle}{Aristotle}{A view of a statue in a Vienna museum.\\
Courtesy of \emph{Scripta Mathematica}.}

\chapter*{\raggedright\Huge LOGIC}
\addcontentsline{toc}{chapter}{2.\ Logic}
\markboth{LOGIC}{}
\vspace{-1em}\noindent\rule{\linewidth}{0.6pt}\vspace{1em}

%========================================================= 2-1
\section*{2--1\quad STATEMENTS}
\addcontentsline{toc}{section}{2--1\quad Statements}

Everyone thinks he knows what a \emph{statement} is. A politician issues a
\emph{statement} when he concedes defeat during an election. The plumber sends
your father a \emph{statement} after installing a new sink. But in mathematics we
use the term in a narrow technical sense. We mentioned in the first chapter that
clear, precise definitions are necessary. Accordingly, we make the following
definition, which turns out to be very useful.

\begin{definition}
A statement is a sentence which is either true or false, but not both.
\end{definition}

Consider the following:

\begin{enumerate}[label=(\alph*), leftmargin=2.2em, itemsep=1pt, topsep=3pt]
\item Go back home.
\item $2 \times 2 = 5$.
\item $2 \times 2 = 4$.
\item New York is a suburb of Los Angeles.
\item Thursday up tar paper.
\item The silliness of speedily.
\end{enumerate}

Some of the above are statements and some are not. Before reading on, decide for
yourself which ones satisfy the definition.

Let us look at some slightly different examples.

\begin{enumerate}[label=(\alph*), leftmargin=2.2em, itemsep=1pt, topsep=3pt]
\item I am the most beautiful girl in town.
\item If I don't get a good night's sleep, I can't pass the examination.
\item $x + 2 = 5$.
\end{enumerate}

All of these are sensible sentences. The only question is whether or not they
satisfy our condition of being either true or false but not both. Before we can
decide whether (a) is a true statement or a false statement we must know who
``I'' is and what town is meant. Then we can decide upon the truth or falsity.
(We could hold a beauty contest!) In (c) if we replace the letter $x$ by ``3,''
we have a true statement. If we replace $x$ by ``2'' the resulting statement is
false. In statement (b), we must again know who ``I'' is and what examination is
meant in order to decide whether it is a true or false statement.

Besides rather simple statements, we can have more complicated ones, such as:

\begin{enumerate}[label=(\alph*), leftmargin=2.2em, itemsep=1pt, topsep=3pt]
\item I cannot go to the dance because I have nothing to wear, my feet hurt, and
      I can't dance.
\item A triangle has three equal sides if and only if it has three equal angles.
\item Jones is a mathematician and Smith is an engineer, while Brown doesn't work
      at all.
\item There is one and only one line through two points.
\item Either I will marry you or get a new job.
\item Two lines either intersect or are parallel.
\end{enumerate}

In some of these examples the full statement is made up of simpler statements. We
see that statements may be combined to give more complicated statements. For
example:

\begin{center}
``Mr.\ Baker is a barber,''\qquad ``Mr.\ Barber is a baker,''
\end{center}

\noindent are two statements. We can combine them into a single statement by
joining them with ``and.''

\begin{center}
``Mr.\ Baker is a barber and Mr.\ Barber is a baker.''
\end{center}

Here are some examples of combining two statements. Make up some others of your
own.

\begin{enumerate}[label=(\alph*), leftmargin=2.2em, itemsep=1pt, topsep=3pt]
\item We are having a picnic and it is not raining.
\item It is night and I can hear the wind.
\item Mary is going with Fred and Suzie is going with John.
\item I like spinach and I dislike candy.
\item You are rich and I am poor.
\item Thursday is April second and the moon is made of green cheese.
\end{enumerate}

To be perfectly clear about combining statements, we indicate statements by
letters, using Greek letters like $\al$ (alpha), $\be$ (beta),
$\ga$ (gamma). Now if we let $\al$ indicate the statement ``Mr.\ Barber is a
baker'' and let $\be$ indicate ``Mr.\ Baker is a barber,'' then the statement
``Mr.\ Barber is a baker and Mr.\ Baker is a barber'' may be written
``$\al$ and $\be$.''

\begin{IIproblem}
In each of the combined statements (a) through (f) above, pick
out the two statements which form it. Note that each statement has the form
``$\al$ and $\be$.''
\end{IIproblem}

The combination of two statements with ``and'' to get a new statement has a name.
We emphasize this by making the definition:

\begin{definition}
If $\al$ and $\be$ are statements, the statement ``$\al$ and $\be$'' is called the
\emph{conjunction} of $\al$ and $\be$.
\end{definition}

You should not accept definitions blindly. It is possible to sneak a false
statement into a definition! Such a definition would be worthless. From the way
Definition 2--2 is worded it is clear that we have \emph{taken for granted} that
``$\al$ and $\be$'' will always be a statement. But a sentence is not a statement
(according to Definition 2--1) unless it is either true or false. Therefore, we
must decide upon some \emph{rule} which will enable us to say when ``$\al$ and
$\be$'' is true and when it is false. If we cannot agree upon such a rule,
Definition 2--2 will be meaningless.

The following statements are all of the form ``$\al$ and $\be$.'' Decide which
ones are true and which ones are false, and then try to state a general rule for
deciding upon the truth of such statements.

\begin{enumerate}[label=(\alph*), leftmargin=2.2em, itemsep=1pt, topsep=3pt]
\item $2 \times 2 = 4$ and $2 + 2 = 5$.
\item $3 + 2 = 5$ and there are at least 52 Sundays in every year.
\item A triangle has three sides and a pentagon has six sides.
\item Every rectangle is a square and all men have purple hair.
\item The positive square root of 16 is 4 and the negative square root of 16 is
      $-4$.
\end{enumerate}

You will probably have no trouble coming to the same conclusion that
mathematicians have agreed upon. That is, ``$\al$ and $\be$'' is considered true
when both $\al$ and $\be$ are true. If one of the statements is false (or both
are), then ``$\al$ and $\be$'' is false. We can describe this rule most clearly by
the \emph{truth table} below.

\begin{truthtable}{ccc}
\tth{\al} & \tth{\be} & \multicolumn{1}{c}{$\al$ and $\be$} \\[2pt]
\TT & \TT & \TT \\
\TT & \FF & \FF \\
\FF & \TT & \FF \\
\FF & \FF & \FF \\
\end{truthtable}

The table is self-explanatory. It points out that ``$\al$ and $\be$'' is true (T)
in only one case and false (F) in all others.

The next step in our development of logic will be harder. We can combine
statements in another way by putting the word ``or'' between them. For example:

\begin{quote}
I will go to the movies \emph{or} I will stay home.\\
I will receive an $A$ in geometry \emph{or} Latin.\\
State University will not admit you unless you have had a year of Spanish
\emph{or} a year of Latin.
\end{quote}

Study the three statements above carefully. Everyone would agree that in the
first statement the speaker means that he will go to the movie \emph{or} he will
stay home \emph{but he will not do both}. In the second statement does the
speaker mean that he will get \emph{either} an $A$ in geometry \emph{or} an $A$
in Latin \emph{but not both}? Or does he mean that he is sure of getting one $A$
and perhaps he will get two? Of course there is no question about the last
statement. State University will not refuse to admit you because you had one year
of Spanish \emph{and} one year of Latin.

We are pointing out that in everyday speech we use the word ``or'' ambiguously.
Sometimes when we say ``$\al$ or $\be$'' we mean that \emph{one and only one} of
the statements is true. Sometimes we mean that \emph{at least} one is true and
possibly both are. We cannot use the word ``or'' with this double meaning in
mathematics. We must explain precisely what we mean when we say ``$\al$ or
$\be$.'' Mathematicians have agreed to call statements of the form ``$\al$ or
$\be$'' true in all cases except when $\al$ and $\be$ are both false. Of course
this differs from the ordinary usage of the word ``or.'' But we use everyday
language too carelessly for our purpose in mathematics. We must be able to
express ourselves precisely and know that our meaning will not be misunderstood.
You must remember that when we say ``$\al$ or $\be$'' we mean that either $\al$
is true and $\be$ is false, or $\al$ is false and $\be$ is true, or $\al$ is true
and $\be$ is true. We make the following formal definition.

\begin{definition}
The statement ``$\al$ or $\be$'' is called the \emph{disjunction} of $\al$ and
$\be$. It is false when both $\al$ and $\be$ are false and true in all other
cases.
\end{definition}

\begin{IIproblem}
Construct the truth table for the statement ``$\al$ or $\be$.''
\end{IIproblem}

\exercisegroup{2--1}

In each exercise below there are two statements. Form their conjunction and
disjunction. Discuss the resulting statements with respect to truth and falsity.

\begin{exercisecols}
\begin{exlist}
\item Some roses are red. Some violets are blue.
\item September has 30 days. There are eight days in a week.
\item 3 is less than 4. 4 is less than 3.
\item Napoleon was an emperor of Japan. Marco Polo discovered America.
\item $2 + 2 = 5$; \quad $6 + 3 = 9$.
\item $7 + 4 = 11$; \quad $8 + 0 = 9$.
\item $1 + 1 = 2$; \quad $1 \times 1 = 2$.
\item $5 + 1 = 5$; \quad $5 \times 1 = 6$.
\end{exlist}
\end{exercisecols}

%========================================================= 2-2
\section*{2--2\quad NEGATION}
\addcontentsline{toc}{section}{2--2\quad Negation}

Sometimes we make statements about other statements. One of the simplest and most
useful statements of this type has the form ``$\al$ is false.'' At times you may
make a statement that you believe true, and someone else may show his
disagreement by saying, ``That is not true.'' Just as we use the terms
\emph{conjunction} and \emph{disjunction} to describe the statements ``$\al$ and
$\be$'' and ``$\al$ or $\be$,'' so also we have a particular name for the
statement ``$\al$ is false.''

\begin{definition}
The statement ``$\al$ is false'' is called the \emph{negation} of $\al$ (and will
be written, and read, ``not $\al$'').
\end{definition}

\begin{IIproblem}
State a rule for determining when ``not $\al$'' is true and when
it is false. Make a truth table for ``not $\al$.''
\end{IIproblem}

When we form the negation of a statement in words we do not usually place ``not''
in front of it. Logically it would be correct to do so, but the resulting
sentence might sound awkward. For example, let $\al$ be ``a square is a
rectangle.'' Anyone untrained in logic would have a right to complain that he did
not understand our meaning if we said, ``not a square is a rectangle.'' The
examples below indicate how negations may be formed. After studying these, make
up several examples of your own.

\begin{enumerate}[label=(\alph*), leftmargin=2.6em, itemsep=3pt, topsep=3pt]
\item $\al$: A square is a rectangle.\\
      not $\al$: A square is not a rectangle.
\item $\al$: $2 + 2 = 4$.\\
      not $\al$: $2 + 2 \neq 4$.
\item $\al$: A triangle has five sides.\\
      not $\al$: A triangle does not have five sides.
\item $\al$: It is false that my donkey has three legs.\\
      not $\al$: My donkey has three legs.
\end{enumerate}

\exercisegroup{2--2}

Form negations of the following statements.

\begin{exercisecols}
\begin{exlist}
\item The grass in my yard is green.
\item This rose is black.
\item I am tired.
\item The temperature of the sun is more than 1000 degrees.
\item Geometry is interesting.
\item $5 + 4 = 9$.
\item $6 + 3 = 7$.
\item[\stex 8.] $x = 2$ and $y = 3$.
\item[\stex 9.] $x = 2$ or $y = 3$.
\item[\stex 10.] All pigs are fat.
\item[\stex 11.] Some Americans are liars.
\end{exlist}
\end{exercisecols}
% p.25 carries the "*" footnote at the foot of the page holding the first
% starred exercises (Ex. Group 2-2, nos. 8-11); anchor it here, not mid-prose.
\IIstarnote{}

Sometimes we wish to state the negation of a statement like ``All cats are
black.'' If we can find \emph{one} cat which is not black we will have proved
this statement false. Hence we form the negation of ``All cats are black'' by
saying ``There is at least one cat which is not black'' or ``Some cat is not
black.'' The most common error made in forming the negation of a
statement like ``All students in this class like geometry'' is to say ``No
student in this class likes geometry.'' These statements certainly do not agree
but they are definitely not negations of each other. Why? Remember that if $\al$
is false then ``not $\al$'' must be true, for ``not $\al$'' states that $\al$ is
false. Show that the statements, ``All students in this class like geometry'' and
``No student in this class likes geometry,'' could both be false.

Another common error is to negate ``All cats are black'' by saying ``All cats are
not black.'' Although in everyday speech this sentence might be understood to
mean ``Not all cats are black,'' we do not interpret it so. If we say ``All cats
are red,'' we mean \emph{every cat has the property of being red}. So we interpret
``All cats are not black'' to mean \emph{every cat has the property of being not
black}. This is a subtle distinction, and again it points out that the precise
language of mathematics differs from careless, everyday speech.

\exercisegroup{2--3}

Give the negation of each statement.

\begin{exercisecols}
\begin{exlist}
\item All men in Georgia are at least six feet tall.
\item All students in France like geometry.
\item All students in Russia hate geometry.
\item Every rectangle is a square.
\item Every triangle is equilateral.
\item All gadgets are widgets.
\item All miggles are maggles.
\item[\stex 8.] For every pair of numbers $a$ and $b$, $a + b = b + a$.
\item[9.] Every circle has a center.
\end{exlist}
\end{exercisecols}

Consider the statement ``\emph{Some} cats are black.'' We usually understand the
word ``some'' to mean ``more than one.'' But in logic it is more convenient to
agree that it means ``one or more.'' (We must have precise definitions!) The
negation of ``Some cats are black'' then becomes ``No cat is black.'' We could
also say, ``Every cat is a color different from black'' or ``Every cat is not
black.'' Notice again that this last statement, ``Every cat is not black,'' means
something altogether different from ``Not every cat is black.''

\exercisegroup{2--4}

Give negations of the following statements. After doing this, make up some
examples of your own.

\begin{exercisecols}
\begin{exlist}
\item Some men play golf.
\item Some boys play football.
\item There is at least one equilateral triangle.
\item For some numbers $x$, $x^2 + 2 = 5$.
\item There is at least one number $x$ such that $x + 3 = 7$.
\item There is at least one number $x$ such that $x^2 + 2 = 0$.
\item Some dogs have fewer than four legs.
\item[\stex 8.] Some numbers are negative.
\item[9.] Some fishermen always tell the truth.
\item[10.] Some people work hard.
\end{exlist}
\end{exercisecols}

When we have two statements, $\al$ and $\be$, and wish to form the \emph{negation}
of their \emph{conjunction}, we must think a little harder than before. Thus if we
wish to form the negation of ``Smith is a dentist and Jones is a teacher,'' we
must say that the statement is false. By our agreement in Section 2--1, we must
say that at least one of the statements ``Smith is a dentist'' and ``Jones is a
teacher'' is false. We do this by saying, ``Smith is \emph{not} a dentist or Jones
is \emph{not} a teacher.''

As another illustration, suppose one student says, ``We will win our game with
Lincoln \emph{and} we will win our game with Bay City.'' To call his statement
false is to say, ``We will \emph{not} win our game with Lincoln \emph{or} we will
\emph{not} win our game with Bay City.'' It should be clear now that the negation
of ``$\al$ and $\be$'' is the statement ``not $\al$ or not $\be$.''

\begin{IIproblem}
Construct a truth table for each of the statements ``$\al$ and
$\be$'' and ``not $\al$ or not $\be$.'' Show that when one statement is true the
other is false, and vice versa.
\end{IIproblem}

\exercisegroup{2--5}

Give the negations of the statements below. Make up some additional exercises of
your own.

\begin{exercisecols}
\begin{exlist}
\item Today is Monday and the weather is cold.
\item $x = 2$ and $y = 5$.
\item $2 \times 2 = 4$ and $3 \times 5 = 15$.
\item $\tri{ABC}$ is isosceles and $\tri{ABC}$ is a right triangle.
\item The point $A$ is on the line $l$ and the point $B$ is on the line $m$.
\item[\stex 6.] All cats are black and all men are mortal.
\item[\stex 7.] All cats are black and some birds are blue.
\item[\stex 8.] Some rectangles are squares and some triangles are equilateral.
\item[\stex 9.] Some widgets are gadgets and no wiggles are waggles.
\end{exlist}
\end{exercisecols}

We may wish to form the \emph{negation} of the \emph{disjunction} of two
statements. For example, we may wish to deny the statement ``Johnson is a tailor
or Brown is a doctor.'' You should have no trouble expressing the negation of this
statement.

\begin{IIproblem}
State the general rule for forming the negation of the
disjunction of two statements. Do so by completing the sentence, ``The negation of
`$\al$ or $\be$' is \underline{\hspace{5em}}.'' Construct a truth table for each of
the statements ``$\al$ or $\be$'' and ``not $\al$ \emph{and} not $\be$.''
\end{IIproblem}

\exercisegroup{2--6}

Give the negation of each statement below. Invent some more examples.

\begin{exercisecols}
\begin{exlist}
\item I will go to the game or I will go to the movie.
\item I will study hard or I will fail.
\item We will start at 6:30 or we will be late.
\item A donkey has three legs or a rabbit has long ears.
\item $2 + 2 = 5$ or $2 + 2 = 6$.
\item $3 \times 7 = 21$ or $4 \times 5 = 20$.
\item Billings is right or Jones is wrong.
\item[\stex 8.] Some will work or all will starve.
\item[\stex 9.] Every triangle is isosceles or every triangle is right.
\item[\stex 10.] No roses are red or some violets are blue.
\end{exlist}
\end{exercisecols}

%========================================================= 2-3
\section*{2--3\quad WHAT ARE PROOFS?}
\addcontentsline{toc}{section}{2--3\quad What are proofs?}

In everyday life when trying to convince a person that something is true, we may
say, ``\emph{This} is true because \emph{that} is true.'' For such a statement to
be convincing, both sides in the argument must agree that the \emph{that} is true
and the \emph{this} must then \emph{necessarily} follow. In other words, there
must be agreement as to what information we start with and how we draw
conclusions from this information. \emph{Logic is the study of the rules for
making correct statements.} Proof always consists of making statements of the
following type: If so and so is true, then such and such must also be true. For
example, consider the statement, ``If the sky is blue (no clouds at all), then it
is not raining.'' Everyone would agree to this statement because from our
experience we know that rain invariably comes from clouds. However, the following
statement, which you might hear in a political speech, is not logically correct.
``The United States is the best country because it is a democracy.'' That we are a
democracy is true; that we are the best country (whatever that means) may be true
although it would not be accepted by all. The logical error comes from the fact
that there are other democratic countries, and hence they too by our careless
reasoning should be best. Therefore, that the United States is best does not
follow from the fact that we are a democracy.

Here are simple examples from algebra:

\begin{center}
``If $3x = 4$, then $x = \tfrac{4}{3}$.''\\[2pt]
``If $a = 2$ and $b = 3$, then $a^2 + ab = 10$.''
\end{center}

Both of these statements are correct and both are of the form, ``If $\al$ then
$\be$.''
\emph{All mathematical proofs employ statements of this type.} We express such a
statement another way by saying

\begin{center}
$\al$ \emph{implies} $\be$.
\end{center}

For example, let $\al$ be the statement ``Today is Monday'' and let $\be$ be the
statement ``Tomorrow is Tuesday.'' Then the statement ``if $\al$ then $\be$'' is
certainly true. However, the statement ``If the sun comes up tomorrow, then it
will not rise the next day'' is very probably false. At least we shall hope that
it turns out to be false. It is fortunate that statements like ``If he takes her
to the dance, then I'll never speak to him again'' and ``If our party loses the
election, then the country will be ruined'' are usually false. It is convenient to
have a symbolic way to indicate that $\al$ implies $\be$. We shall write this as

\begin{center}
$\al \imp \be$,
\end{center}

\noindent where you should read the arrow as the word ``implies.'' We call the
statement ``$\al$ implies $\be$'' an \emph{implication}.

\exercisegroup{2--7}

Put the following statements in an ``if $\al$ then $\be$'' form so that $\al$
implies $\be$.

\begin{exercisecols}
\begin{exlist}
\item If I eat strawberries I get hives.
\item The picnic will be canceled in case of rain.
\item We will win the game if Jonesevitch plays.
\item Numbers larger than 5 are positive.
\item $x^2 = y^2$ because $x = y$.
\item $3x^2 - 2 = 10$ when $x = 2$.
\item A geometric figure which is a square is not a circle.
\item Squares are rectangles.
\item The area of a rectangle with sides of length $a$ and $b$ is $ab$.
\item With a million dollars, I'd eat steak for breakfast.
\item People who don't study mathematics don't have good sense.
\item When it rains my rheumatism bothers me.
\item When I study I do well in school.
\item People who do not vote are poor citizens.
\item Two lines which intersect are not parallel.
\item I know that $x = 2$ because $3x = 6$.
\item Satisfaction guaranteed or your money back.
\item I know $3x = 6$ because $x = 2$.
\item I will not go to the party if Jimmy is coming.
\item This shirt will not shrink or your money will be refunded.
\item This shirt will shrink or your money will be refunded.
\item To remain free we must vote in elections.
\item The crops will be saved if we have rain this week.
\item The crops will be saved only if we have rain this week.
\item I will go to the dance if Joe asks me.
\item I will go to the dance only if Joe asks me.
\item It will rain tomorrow if it turns cold.
\item It will rain tomorrow only if it turns cold.
\end{exlist}
\end{exercisecols}

Working with statements of the form ``$\al \imp \be$'' is called \emph{deductive
reasoning} or, in everyday language, \emph{drawing conclusions}. The statement
$\al$ is called the \emph{hypothesis} and the statement $\be$ is called the
\emph{conclusion}. We use this method of reasoning in daily life when we try to
prove to someone that some statement $\be$ is true. We try to find another
statement $\al$ which implies $\be$, and which is true. Then we say, ``$\al$ is
true, therefore $\be$ is true.''

\exercisegroup{2--8}

In Exercises 1 through 7, a statement $\al \imp \be$ is given which is to be
accepted as true. What conclusion can you draw if (a) is true? if (b) is true?

\begin{exercisecols}
\begin{exlist}
\item If Fred is on the team he must go to bed early.\\
      (a) Fred goes to bed early.\qquad (b) Fred does not go to bed early.
\item If a man is a postman he has foot trouble.\\
      (a) This man has perfect feet.\qquad (b) That man has foot trouble.
\item You will be given a ticket if you are a sophomore.\\
      (a) Jim was not given a ticket.\qquad (b) Joe was given a ticket.
\item You will be given a ticket only if you are a sophomore.\\
      (a) Jean was given a ticket.\qquad (b) George is a sophomore.
\item If two triangles have equal bases and heights then they have equal areas.\\
      (a) These triangles have equal areas.\qquad (b) These triangles do not have
      equal areas and they do have equal heights.
\item The house will burn down if we don't get some water.\\
      (a) We did not get any water.\qquad (b) We did get some water.
\item If a man is beheaded in June then he will die before July 4.\\
      (a) This man died before July 4.\qquad (b) Smith was not beheaded in June.
\end{exlist}

Is the reasoning in Exercises 8 through 18 correct? Invent some other examples
illustrating both correct and incorrect reasoning.

\begin{exlist}\setcounter{exlisti}{7}
\item All men are mortal. Socrates is a man.\\
      \emph{Conclusion:} Socrates is mortal.
\item All widgets are gadgets. This is a gadget.\\
      \emph{Conclusion:} This is a widget.
\item All widgets are gadgets. This is not a gadget.\\
      \emph{Conclusion:} This is not a widget.
\item Only widgets are gadgets. This is a gadget.\\
      \emph{Conclusion:} This is a widget.
\item Only widgets are gadgets. This is a widget.\\
      \emph{Conclusion:} This is a gadget.
\item Every equilateral triangle is isosceles. $\tri{ABC}$ is isosceles.\\
      \emph{Conclusion:} $\tri{ABC}$ is equilateral.
\item Parallel lines do not intersect. Lines $r$ and $s$ intersect.\\
      \emph{Conclusion:} Lines $r$ and $s$ are not parallel.
\item All of these dresses will shrink if washed in hot water. This dress did
      shrink.\\
      \emph{Conclusion:} It was washed in hot water.
\item All sweet food has honey in it. This food is sweet.\\
      \emph{Conclusion:} It has honey in it.
\item All sweet food has honey in it. This food doesn't have honey in it.\\
      \emph{Conclusion:} It is not sweet.
\item All sweet food has honey in it. This food is not sweet.\\
      \emph{Conclusion:} It does not have honey in it.
\end{exlist}
\end{exercisecols}

%========================================================= 2-4
\section*{2--4\quad THE TRUTH TABLE FOR IMPLICATION}
\addcontentsline{toc}{section}{2--4\quad The truth table for implication}

If $\al$ and $\be$ are statements, then ``$\al \imp \be$'' should be another
statement. In other words, ``$\al \imp \be$'' must be either true or false. In
that case we must agree on a way of determining whether ``$\al \imp \be$'' is true
or false.

The definition we give may seem a bit strange at first, but keep in mind two
things that may help you. In the first place the definition below is precise. In
the second place we want to use the implication $\al \imp \be$ only when the
hypothesis $\al$ is presumed true.

\begin{definition}
$\al \imp \be$ is true whenever the statement ``(not $\al$) or $\be$'' is true. It
is false when the statement ``(not $\al$) or $\be$'' is false.
\end{definition}

The truth table, then, is as follows.

\begin{truthtable}{cccc}
\tth{\al} & \tth{\be} & \multicolumn{1}{c}{$\al \imp \be$}
          & \multicolumn{1}{c}{(not $\al$) or $\be$} \\[2pt]
\TT & \TT & \TT & \TT \\
\TT & \FF & \FF & \FF \\
\FF & \FF & \TT & \TT \\
\FF & \TT & \TT & \TT \\
\end{truthtable}

The last line may surprise you. Just remember that $\al \imp \be$ means that
$\be$ is true or $\al$ is false and that $\be$ is true \emph{if} $\al$ is true.

\exercisegroup{2--9}

The exercises referred to below are those in Exercise Group 2--8.

\begin{exercisecols}
\begin{exlist}
\item Write each statement $\al \imp \be$ in Exercises 1 through 7 in the
      ``(not $\al$) or $\be$'' form.
\item Re-examine the validity of the conclusions in Exercises 8 through 18.
\item A statement $\al \imp \be$ is also sometimes expressed as ``$\al$ is
      sufficient for $\be$,'' or ``$\al$ suffices for $\be$,'' or again, ``for
      $\be$ to be true it is sufficient that $\al$ be true.'' Thus for $x^2$ to
      equal 4 it is sufficient that $x = -2$. Express the implications in
      Exercises 1 through 7 in this language.
\item A statement $\al \imp \be$ is also sometimes expressed as ``$\be$ is
      necessary for $\al$'' or ``in order that $\al$ be true it is necessary that
      $\be$ be true.'' Thus for $x$ to equal $-2$ it is necessary that $x^2 = 4$.
      Express the implications in the Exercises 1 through 7 in this language.
\end{exlist}
\end{exercisecols}


%========================================================= 2-5
\section*{2--5\quad THE CONVERSE OF AN IMPLICATION}
\addcontentsline{toc}{section}{2--5\quad The converse of an implication}

From one implication ``$\al \imp \be$'' we can form others which may or may not be
true. One way of doing this gives what we call the \emph{converse}.

\begin{definition}
The \emph{converse} of ``$\al \imp \be$'' is ``$\be \imp \al$.''
\end{definition}

\emph{Example 1.} $\al \imp \be$: If an animal is a man then it is a human being.\\
\emph{Converse:} $\be \imp \al$: If an animal is a human being then it is a man.\\
In this example ``$\al \imp \be$'' is true but ``$\be \imp \al$'' is false.

\medskip
\emph{Example 2.} $\al \imp \be$: If $\tri{ABC}$ is equilateral then $\tri{ABC}$
is equiangular.\\
\emph{Converse:} $\be \imp \al$: If $\tri{ABC}$ is equiangular, then $\tri{ABC}$
is equilateral.\\
In this example both ``$\al \imp \be$'' and ``$\be \imp \al$'' are true.

\exercisegroup{2--10}

Write the converses of the following statements. In each exercise decide whether
the statement is true or false. Do the same thing for the converse.

\begin{exercisecols}
\begin{exlist}
\item Every man born in Indiana is an American citizen.
\item If $x - y = 3$, then $x$ is greater than $y$.
\item If two rectangles have equal bases and equal altitudes, they have equal
      areas.
\item If $x^2 = 9$, then $x = 3$ or $x = -3$.
\item If a man is a Polynesian, he is at least six feet tall.
\item If plants are green, they must have had sunshine.
\item If a point is on the perpendicular bisector of a line segment, it is
      equidistant from the endpoints of the segment.
\item If I am sleeping, then I am breathing.
\item If $2 + 4 = 5$, then $4 + 7 = 10$.
\item If $2 + 4 = 5$, then $4 + 7 = 11$.
\item If $2 + 4 = 6$, then $4 + 7 = 10$.
\item If $2 + 4 = 6$, then $4 + 7 = 11$.
\end{exlist}
\end{exercisecols}

%========================================================= 2-6
\section*{2--6\quad EQUIVALENT STATEMENTS}
\addcontentsline{toc}{section}{2--6\quad Equivalent statements}

When two people give the same mathematical proof they probably will not make
exactly the same statements. But it is very likely that at certain places in their
proofs they will make statements which we say are \emph{equivalent} to one
another. In order to understand precisely what this means we make the following
definition.

\begin{definition}
Two statements $\al$, $\be$ are said to be \emph{equivalent} to one another if
\emph{both} the statement $\al \imp \be$ and the converse $\be \imp \al$ are true.
If $\al$ and $\be$ are equivalent statements, we indicate this by writing
$\al \iffx \be$, which we read as ``$\al$ implies $\be$ and $\be$ implies $\al$.''
\end{definition}

This definition may be difficult to understand. An explanation is that
``equivalent statements present the same information---they may simply be stated
differently.''

An example of a pair of equivalent statements is

\begin{quote}
$\al$: Line $l$ is perpendicular to line $m$.\\
$\be$: Line $m$ is perpendicular to line $l$.
\end{quote}

\exercisegroup{2--11}

In 1 through 7 below decide whether the statements are equivalent to one another.
Invent a few examples of pairs of statements which are equivalent to one another
and some which are not.

\begin{exercisecols}
\begin{exlist}
\item $\al$: Point $P$ lies on the circle with center at $O$ and radius five
      inches.\\
      $\be$: The length of the line segment $OP$ is five inches.
\item $\al$: Line $l$ is parallel to line $m$.\\
      $\be$: Line $n$ is parallel to line $m$.
\item $\al$: $x$ is greater than $y$.\\
      $\be$: $y$ is less than $x$.
\item $\al$: $2x + 3 = 17$.\\
      $\be$: $x = 7$.
\item $\al$: This geometric figure is a square.\\
      $\be$: This geometric figure is a special kind of parallelogram.
\item $\al$: $\tri{ABC}$ is equilateral.\\
      $\be$: $\tri{ABC}$ is equiangular.
\item $\al$: $x = 3$ or $x = -3$.\\
      $\be$: $x^2 = 9$.
\item[\stex 8.] A statement $\al \iffx \be$ is also sometimes expressed by
      saying, ``In order that $\al$ be true it is necessary and sufficient that
      $\be$ be true.'' In each of the above problems state the equivalence of
      $\al$ and $\be$ in this language.
\end{exlist}
\end{exercisecols}

Exercises 4 and 7 above illustrate that equivalent statements have an important
place in algebra. See if you can prove that the following statement is true:

\begin{quote}
\emph{If two statements $\al$, $\be$ are equivalent to each other, then either
both are true or both are false.}
\end{quote}

An easy way to prove this is to complete the following truth table. Can you
understand how this table really presents a proof?

\begin{truthtable}{ccccc}
\tth{\al} & \tth{\be} & \multicolumn{1}{c}{$\al \imp \be$}
          & \multicolumn{1}{c}{$\be \imp \al$}
          & \multicolumn{1}{c}{$\al \iffx \be$} \\[2pt]
\TT & \TT & \TT & \TT & ? \\
\TT & \FF & \FF & \TT & ? \\
\FF & \TT & \TT & \FF & ? \\
\FF & \FF & \TT & \TT & ? \\
\end{truthtable}

This truth table must be used together with Definition 2--7. By this definition we
know whether to mark $\al \iffx \be$ as true or false. Studying the table, we see
that according to Definition 2--7, the statement $\al \iffx \be$ must be labeled
true on lines 1 and 4 and labeled false on lines 2 and 3. It should be easy for
you to complete the proof.

\exercisegroup{2--12}

Form statements equivalent to those listed below. If you are not sure of the
equivalence, make a truth table and check your work.

\begin{exercisecols}
\begin{exlist}
\item $x + y = 7$ and $x - y = 3$.
\item $x$ is greater than $y$.
\item Line segments $AB$ and $CD$ do not intersect.
\item There is a point $P$ which lies on both of the lines $l$ and $m$.
\item Whether you are a smoker or whether you are not, whether you are married or
      single, whether you are a church member or not, you will like Dum-Dums.
\item This geometric figure has four sides, and pairs of opposite sides are
      parallel to one another.
\item That figure is a triangle.
\item This geometric figure consists of two points $A$ and $B$ and all the points
      between $A$ and $B$.
\item The points $A$ and $B$ lie on the same side of the line $l$.
% p.33 prints "opposite lines of the line l" -- kept verbatim; see UNCERTAIN.
\item The points $A$ and $B$ lie on opposite lines of the line $l$.
\end{exlist}
\end{exercisecols}

%========================================================= 2-7
\section*{2--7\quad THE CONTRAPOSITIVE OF AN IMPLICATION}
\addcontentsline{toc}{section}{2--7\quad The contrapositive of an implication}

A second important implication related to $\al \imp \be$ is its
\emph{contrapositive}.

\begin{definition}
The statement ``not $\be \imp$ not $\al$'' is called the \emph{contrapositive} of
the statement ``$\al \imp \be$.''
\end{definition}

\emph{Examples:}

\begin{enumerate}[label=(\alph*), leftmargin=2.6em, itemsep=3pt, topsep=3pt]
\item $\al \imp \be$: $x = 3 \imp x^2 = 9$.\\
      not $\be \imp$ not $\al$: $x^2 \neq 9 \imp x \neq 3$.
\item $\al \imp \be$: If two lines intersect they are not parallel.\\
      not $\be \imp$ not $\al$: If two lines are parallel they do not intersect.
\item $\al \imp \be$: If I make an $A$ on the examination I will receive an $A$ in
      the course.\\
      not $\be \imp$ not $\al$: If I do not receive an $A$ in the course then I
      will not have made an $A$ on the examination.
\item $\al \imp \be$: If I go to the game I will not go to the movie.\\
      not $\be \imp$ not $\al$: If I go to the movie I will not go to the game.
\end{enumerate}

The most interesting property of the statement ``not $\be \imp$ not $\al$'' is
that it is equivalent to $\al \imp \be$. We can prove this by showing that these
two implications have identical truth tables.

\begin{IIproblem}
Complete the following truth table and show that ``$\al \imp
\be$'' and ``not $\be \imp$ not $\al$'' are equivalent to each other.
\end{IIproblem}

\begin{truthtable}{cccccc}
\tth{\al} & \tth{\be} & \multicolumn{1}{c}{not $\al$}
          & \multicolumn{1}{c}{not $\be$}
          & \multicolumn{1}{c}{$\al \imp \be$}
          & \multicolumn{1}{c}{not $\be \imp$ not $\al$} \\[2pt]
\TT & \TT & \FF & \FF & --- & --- \\
\TT & \FF & \FF & \TT & --- & --- \\
\FF & \TT & \TT & \FF & --- & --- \\
\FF & \FF & \TT & \TT & --- & --- \\
\end{truthtable}

\exercisegroup{2--13}

Form the contrapositive of each statement in 1 through 8 below.

\begin{exercisecols}
\begin{exlist}
\item If $2 \times 2 = 4$, then oranges are composed of uranium.
\item If you can run the mile in four minutes, then I'll eat my hat.
\item If the ground hog sees his shadow, then we will have a cold spring.
\item If there is no atmosphere on Mars, then there is no life there.
\item If point $P$ is on line $r$, then it is not on line $s$.
\item If you are not satisfied, then your money will be returned.
\item If food is sweet, then it has honey in it.
\item If you always obey traffic laws, then you will not be given a ticket.
\item[\stex 9.] Using the word \emph{necessary}, state the contrapositive of the
      implication of Exercise 4.
\item[\stex 10.] Using the word \emph{sufficient}, state the contrapositive of the
      implication of Exercise 4.
\end{exlist}
\end{exercisecols}

%========================================================= 2-8
\section*{2--8\quad ``IF AND ONLY IF''}
\addcontentsline{toc}{section}{2--8\quad ``If and only if''}

Consider the two statements ``$\al$ if $\be$'' and ``$\al$ only if $\be$.'' Our
experience in working with statements like these tells us that ``$\al$ if $\be$''
means ``$\be \imp \al$'' and that ``$\al$ only if $\be$'' means ``$\al \imp \be$.''
If both of these statements are true, we have ``$\al \iffx \be$.'' If both of
these statements are true, their conjunction is true.

\begin{definition}
``$\al$ if and only if $\be$'' means ``$\al \imp \be$ and $\be \imp \al$,'' that
is, ``$\al \iffx \be$.''
\end{definition}

Consider the following two statements. ``Two lines are parallel if they do not
intersect.'' ``Two lines do not intersect if they are parallel.'' According to
Definition 2--9 we may write the conjunction of these two statements as ``Two
lines are parallel if and only if they do not intersect.''

\exercisegroup{2--14}

In each exercise below there are two statements. Formulate a single statement,
using ``if and only if,'' that has the same meaning as the conjunction of these
statements.

\begin{exercisecols}
\begin{exlist}
\item If the weather is bad we will stay home. If we stay home the weather is bad.
\item If you are a sophomore then you will get a ticket. If you get a ticket then
      you are a sophomore.
\item If the weather is bad we will stay home. If the weather is not bad we will
      not stay home.
\item Everyone who receives more than 50\% on this test will pass the course. All
      others will fail.
\item Your dog will obey you if you treat him well. If you do not treat him well
      he will not obey you.
\item I will marry you if you buy me a mink coat. If you don't buy me a mink coat
      I won't marry you.
\item If food is sweet then it has honey in it. Food with honey in it is sweet.
\end{exlist}
\end{exercisecols}

%========================================================= 2-9
\section*{2--9\quad THEOREMS AND PROOFS}
\addcontentsline{toc}{section}{2--9\quad Theorems and proofs}

In Section 2--3 we said logic is a study of the rules for making correct
statements. The Greek philosopher Aristotle was the first great student of logic.
His work was not improved upon until modern times. In fact, all the logical
concepts we have studied in this chapter belong to what we call
\emph{Aristotelian logic}, which is over 2000 years old and is just as useful
today as it ever was. The logical principles of this chapter are all that you will
need for many years. If you have found these ideas interesting, you may wish to
take courses in logic later on in college. Whether or not you continue your study
of logic, the work of this chapter should help you become careful in both speech
and writing. In particular, during the work in geometry this year, do not tolerate
false statements. Remember that if one false statement is accepted as true then
strange things can be proved. Let us strive then to express ourselves precisely.

Most of our geometrical work will consist of \emph{proving theorems}. Theorems are
statements. To \emph{prove a theorem} is to show that it follows logically from
our assumptions. To do so we employ our principles of logic. For example, \emph{if}
we know that the statement $\al \imp \be$ is true \emph{and if} we know that $\al$
is true, \emph{then} we can infer that $\be$ is true. In symbols,

\begin{center}
$\al$ and $(\al \imp \be) \imp \be$.
\end{center}

Another technique of proof is the following: \emph{If} we know that $\al \imp \be$
is true \emph{and if} we know that $\be$ is false, \emph{then} we can infer that
$\al$ is false. Symbolically,

\begin{center}
not $\be$ and $(\al \imp \be) \imp$ not $\al$.
\end{center}

You need only look at the truth table for $\al \imp \be$ to realize that the two
rules above are correct.

\begin{truthtable}{ccc}
\tth{\al} & \tth{\be} & \multicolumn{1}{c}{$\al \imp \be$} \\[2pt]
\TT & \TT & \TT \\
\TT & \FF & \FF \\
\FF & \TT & \TT \\
\FF & \FF & \TT \\
\end{truthtable}

\begin{IIproblem}
Explain how this truth table justifies the two methods of proof
described above.
\end{IIproblem}

Let us illustrate how these proof techniques may be employed (although you have
already reasoned in this way in many exercises in this chapter).

\medskip
\emph{Example 1.} The notice said, ``If it rains on commencement day then the
graduation ceremony will be held in the gymnasium.'' ($\al \imp \be$.) Today is
commencement day and it is raining. ($\al$ is true.) Therefore, since both of the
statements $\al$ and $\al \imp \be$ are true, we know that $\be$ is true. That is,
the graduation ceremony will be held in the gymnasium.

\medskip
\emph{Example 2.} Jim says that $-3$ is a root of the equation $x^2 - x - 6 = 0$.
If $-3$ is a root of the equation $x^2 - x - 6 = 0$, then $(-3)^2 - (-3) - 6 = 0$.
($\al \imp \be$.) But, $(-3)^2 - (-3) - 6 = 6$. ($\be$ is false.) Now, since
$\al \imp \be$ is true and $\be$ is false, I know that $\al$ is false and Jim was
wrong. The number $-3$ is not a root of the equation $x^2 - x - 6 = 0$.

\medskip
In these examples we would say that we have proved the theorems,

\begin{exlist}
\item The graduation ceremony will be held in the gym.
\item The number $-3$ is not a root of $x^2 - x - 6 = 0$.
\end{exlist}

We might have written out our proofs as follows, in a ``statement-reason'' form.

\begin{steps}
\step{Rain on commencement day $\imp$ ceremony in gym.}{Given.}
\step{It is raining on commencement day.}{Given.}
\step{The ceremony will be in the gym.}{Statements 1 and 2 and laws of logic.}
\end{steps}

\begin{steps}
\step{$-3$ is a root of $x^2 - x - 6 = 0 \imp (-3)^2 - (-3) - 6 = 0$.}{Definition
      of root.}
\step{$(-3)^2 - (-3) - 6 \neq 0$.}{Arithmetic.}
\step{$-3$ is not a root of $x^2 - x - 6 = 0$.}{Statements 1 and 2 and laws of
      logic.}
\end{steps}

\exercisegroup{2--15}

In each problem below you are given certain statements which we shall call true.
Then a theorem is stated. In some cases the theorem may be false; in others there
may not be enough information given to decide whether the theorem is true or
false. Write out proofs. We refer to our given true statements in each problem as
our \emph{hypothesis} (hyp).

\begin{exercisecols}
\begin{exlist}
\item \begin{hyp}Uncle George will take us skating today if Mother says we may go.
      Mother has said we may go.\end{hyp}
      \begin{thmclaim}Uncle George will take us skating today.\end{thmclaim}
\item \begin{hyp}If there is no oxygen upon the moon then there can be no animal
      life there. Tests have shown conclusively that there is no oxygen upon the
      moon.\end{hyp}
      \begin{thmclaim}There is no animal life on the moon.\end{thmclaim}
\item \begin{hyp}People with green eyes cannot be trusted. Joan has green
      eyes.\end{hyp}
      \begin{thmclaim}Joan cannot be trusted.\end{thmclaim}
\item \begin{hyp}People with blue eyes are honest. Henry is dishonest.\end{hyp}
      \begin{thmclaim}Henry does not have blue eyes.\end{thmclaim}
\item \begin{hyp}If Joe parts his hair on the left then he does not like spinach.
      Joe does not like spinach.\end{hyp}
      \begin{thmclaim}Joe parts his hair on the left.\end{thmclaim}
\item \begin{hyp}Cats with green whiskers have long tails. Cats who do not like
      milk do not have long tails. Simba is a cat who does not like milk.\end{hyp}
      \begin{thmclaim}Simba does not have green whiskers.\end{thmclaim}
\item \begin{hyp}A student who attends M.I.T. will spend more than \$1200 during
      the school year. Kate spent \$980 last year in school.\end{hyp}
      \begin{thmclaim}Kate did not attend M.I.T. last year.\end{thmclaim}
\item \begin{hyp}$x + y = 20$. $x - y = 4$.\end{hyp}
      \begin{thmclaim}$x \neq 11$.\end{thmclaim}
\item \begin{hyp}$2x - 3y = 7$. $x + 2y = 3$.\end{hyp}
      \begin{thmclaim}$3x - y = 10$.\end{thmclaim}
\item \begin{hyp}Joe has 11 marbles: four blue ones and seven yellow ones. Joe
      gave his sister six marbles. Joe did not give all his blue marbles
      away.\end{hyp}
      \begin{thmclaim}Joe does not have more than four yellow marbles
      left.\end{thmclaim}
\item \begin{hyp}Only those students whose test grades average more than 90\% will
      receive $A$'s for their semester's work in geometry. Jim's test grades have
      been 89, 84, 92, 85, and 100. Sally's scores are 91, 90, 89, 94, and
      87.\end{hyp}
      \begin{thmclaim}Sally will receive an $A$ and Jim will not.\end{thmclaim}
\item \begin{hyp}The lovely actress Amanda Thistle has promised to marry her
      honest admirer Silas Moneybags only if he buys her a mink stole, a yacht,
      and the great Calcutta Emerald. Silas has purchased for his Amanda these
      items and many others. Amanda always keeps her promises.\end{hyp}
      \begin{thmclaim}Amanda will wed Silas.\end{thmclaim}
\item \begin{hyp}$x + y = 10$. $x$ and $y$ are positive whole numbers. $x$ is
      greater than $y$. $y + 4$ is greater than $x$.\end{hyp}
      \begin{thmclaim}$y = 4$.\end{thmclaim}
\item \begin{hyp}The square of a positive number is positive. The square of a
      negative number is positive. The square of zero is zero. Zero is neither
      positive nor negative. $k$ is a number. $k^2 = -1$.\end{hyp}
      \begin{thmclaim}There are numbers which are not positive or negative or
      equal to zero.\end{thmclaim}
\item \begin{hyp}The product of two positive numbers is positive. The number $a$
      is positive. The product of numbers $a$ and $b$ is not positive.\end{hyp}
      \begin{thmclaim}The number $b$ is negative.\end{thmclaim}
\item[\stex 16.] \begin{hyp}If $a$ is any number, then $1 \cdot a = a$. If $a$,
      $b$, $c$ are any numbers, then $ba + ca = (b + c)a$. $1 + 1 = 2$.\end{hyp}
      \begin{thmclaim}For every number $a$, it is true that $a + a = 2 \cdot
      a$.\end{thmclaim}
\item[\stex 17.] \begin{hyp}$4 + (-4) = 0$. $6 = 2 + 4$. For all numbers $a$, $b$,
      $c$, it is true that $(a + b) + c = a + (b + c)$. We always have $x + 0 =
      x$.\end{hyp}
      \begin{thmclaim}$6 + (-4) = 2$.\end{thmclaim}
\item[\stex 18.] \begin{hyp}If $8 + x = 0$ then $x = -8$. $8 = 3 + 5$. Numbers can
      be rearranged when we add them. $3 + (-3) = 0$. $5 + (-5) = 0$. $0 + 0 =
      0$.\end{hyp}
      \begin{thmclaim}$(-3) + (-5) = -8$.\end{thmclaim}
\item[\stex 19.] \begin{hyp}If neither of two numbers is equal to zero then their
      product is not zero. The product of the two numbers $a$ and $b$ is zero.
      $b \neq 0$.\end{hyp}
      \begin{thmclaim}$a = 0$.\end{thmclaim}
\item[\stex 20.] \begin{hyp}Dogs with long ears have short tails. Dogs who chase
      rabbits never get measles. Dogs who do not chase rabbits have long tails. My
      dog Rover died of the measles.\end{hyp}
      \begin{thmclaim}Rover did not have long ears.\end{thmclaim}
\end{exlist}
\end{exercisecols}

%========================================================= review
\section*{REVIEW OF CHAPTER 2}
\addcontentsline{toc}{section}{Review of Chapter 2}

The ideas in this chapter are not to be memorized as you may have memorized the
addition and multiplication tables when you studied arithmetic. It is very
important that you realize that \emph{all} our complicated mathematical statements
may be built up from simple statements using only the connectives ``and,'' ``or,''
``not,'' and ``if \dots\ then.'' You should relate the precise meaning of other
connectives like ``only if'' and ``if and only if'' to these four fundamental
terms. \emph{Do not} make any effort to memorize the meanings of these latter
terms. You should remember the meaning of conjunction, disjunction, negation,
implication, and equivalence. You should be able to form the converse or
contrapositive of an implication. You should be able to give examples showing that
two statements, ``$\al \imp \be$,'' and the converse, ``$\be \imp \al$,'' are not
necessarily equivalent to each other, that is, that one might be true and the
other false. You should see clearly that the implication ``$\al \imp \be$'' and its
contrapositive, ``not $\be \imp$ not $\al$,'' really assert the same thing; hence
they are equivalent statements. You should understand and be able to use the
important rules of logic employed so often in proofs, namely,

\begin{exlist}
\item $\al$ and $(\al \imp \be) \imp \be$.
\item not $\be$ and $(\al \imp \be) \imp$ not $\al$.
\end{exlist}

You should understand how to set up and use truth tables. For example, to set up a
truth table for the statement ``$\al$ and $\al \imp \be$,'' we could list the
following:

\begin{truthtable}{cccc}
\tth{\al} & \tth{\be} & \multicolumn{1}{c}{$\al \imp \be$}
          & \multicolumn{1}{c}{$\al$ and $(\al \imp \be)$} \\[2pt]
\TT & \TT & \TT & ? \\
\TT & \FF & \FF & ? \\
\FF & \TT & \TT & ? \\
\FF & \FF & \TT & ? \\
\end{truthtable}

\begin{IIproblem}
Fill in the last column in the truth table above and show that
when the statement ``$\al$ and $\al \imp \be$'' is true, then $\be$ is true also.
Make a similar truth table to illustrate the second rule of logic listed above.
Can you invent a rule of logic like these two and show that it is correct by
setting up a truth table?
\end{IIproblem}

\subsection*{Review Exercises}

\begin{exlist}
\item Under what condition is ``$\al$ and $\be$'' a true statement? Give examples.
\item Give an example to show why we call ``$\al$ or $\be$'' true when both $\al$
      and $\be$ are true.
\item What is the negation of not $\al$? Give examples.
\item Make up a statement of the form ``$\al$ and $\be$.'' State its negation.
\item Make up a statement of the form ``$\al$ or $\be$.'' State its negation.
\item What is meant by ``$\al \imp \be$''? Give two statements having this form,
      one of which is true and the other false.
\item Why is the statement ``$\al$ or not $\al$'' true no matter whether $\al$ is
      true or false? Give examples to illustrate this.
\item[\stex 8.]\stepcounter{exlisti} Construct a truth table to show under what condition
      ``($\al$ and $\be$) or $\ga$'' is a true statement. Invent an example of
      this type.
\item[9.]\stepcounter{exlisti} Give an example of an implication ``$\al \imp \be$'' whose converse is
      true and one whose converse is false.
\item Form the converse of:\\
      (a) If a man is from Chicago then he is rich.\\
      (b) If she sings in the choir then I will not go to church.\\
      (c) If Smith is elected president he will get ulcers.\\
      (d) You must be a sophomore if you are taking geometry.\\
      (e) Only the brave are free.
\item Form the contrapositive of each implication in Exercise 10.
\item Write a statement equivalent to:\\
      (a) Sweet food always has honey in it.\\
      (b) My rheumatism never bothers me in warm weather.\\
      (c) In order to be healthy you must drink milk.\\
      (d) A grade of 80\% on the final examination will be sufficient to pass the
      course.\\
      (e) Sophomores may come to the dance, but no one else will be permitted to
      do so.
\item Write the negation of:\\
      (a) Honesty is the best policy.\\
      (b) All math teachers are brutes.\\
      (c) The potato crop will fail and everyone will starve.\\
      (d) Roses are red and violets are blue, sugar is sweet and I love you.\\
      (e) Some cats can swim but no chickens can.
\end{exlist}

There are two statements in each exercise below. Calling the first statement $\al$
and the second one $\be$, rewrite them in the symbolic form ``$\al$ and $\be$,''
``$\al$ or $\be$,'' ``$\al \imp \be$,'' ``$\be \imp \al$,'' ``$\al \imp$ not
$\be$,'' ``$\al \iffx \be$,'' etc., as the case may be. For example, the statement,
``A man is honest if he has big feet'' would be written ``$\be \imp \al$,'' not
``$\al \imp \be$.''

\begin{exlist}\setcounter{exlisti}{13}
\item Vacation begins May 12 but we are not going any place.
\item A man on the faculty must park his car in the main parking lot.
\item You will not receive an $A$ in geometry if you do not work this example
      correctly.
\item Mother will let me go to the dance if and only if I wash the dishes for a
      week.
\item Students who laugh at his jokes are the only ones who get good grades.
\item It is necessary that we have rain this week if the melon crop is to be good.
\item You will write a term paper or you will fail.
\item This is the last exercise and I am glad that it is.
\end{exlist}

\subsection*{ALGEBRA REVIEW}

In the exercises below we indicate some of the ways that logical reasoning is used
in algebra. First we shall list some statements that we agree to call \emph{true}
in algebra. Then we shall set up algebraic problems which provide opportunities to
use these true statements and reason logically to correct conclusions. In what
follows, all letters refer to integers, that is, whole numbers in the set
$(\dots, -3, -2, -1, 0, 1, 2, \dots)$. Let us agree that the following statements
are true:

\begin{exlist}
\item For all integers $a$, $b$, $c$ it is true that if $a + b = a + c$, then
      $b = c$.
\item For all integers $a$, $b$, $c$ it is true that if $ab = ac$ and $a \neq 0$,
      then $b = c$.
\end{exlist}

Let us prove the theorem: \emph{If $x = 5$ then $3x + 4 = 19$.}

\medskip
\noindent\emph{Proof.} If $x = 5$ then $3x + 4 = 3 \cdot 5 + 4 = 15 + 4 = 19$.

\medskip
Let us prove the converse of the above theorem, namely: \emph{If $3x + 4 = 19$
then $x = 5$.}

\medskip
\noindent\emph{Proof.} $3x + 4 = 19 \imp 3x + 4 = 15 + 4$. Now, by statement (1)
above we know that $3x + 4 = 15 + 4 \imp 3x = 15$ and that $3x = 15 \imp 3x =
3 \cdot 5$. Using statement (2) above, we have $3x = 3 \cdot 5 \imp x = 5$. We
could present this proof briefly in the form:

\begin{center}
$3x + 4 = 19 \imp 3x + 4 = 3 \cdot 5 + 4 \imp 3x = 3 \cdot 5 \imp x = 5$.
\end{center}

\subsection*{Exercises}

\begin{multicols}{2}
\begin{exlist}
\item Prove that if $x = 7$ then $4x + 5 = 33$.
\item Prove the converse of the statement in Exercise 1.
\item Prove that if $7x + 9 = 65$ then $x = 8$, and conversely.
\item Prove that $5x + (-3) = 27$ only if $x = 6$.
\item Prove that $x = 11$ is a necessary condition for $5x + 4 = 59$.
\item Prove that if $x = 11$, $5x + 4 = 59$.
\item Prove that $2x + 7 = 19$ if and only if $4x - 3 = 21$.
\item Prove that there is no number $x$ such that $3x + 5 = 11$ and $4x - 1 = 3$.
\item What can you infer if it is known that $a \neq 0$ and $a(3b - 8) = a$?
\item What can you prove if it is known that $r(7b - 5) = r$?
\item It is given that $(3a + 12)(2b - 8) = 0$ and $b \neq 4$. What can be proved?
\item[\stex 12.] It is known that $(2a - 3)(b - 3) = 2a - 2$. Prove that if
      $a = 1$ then $b = 3$. Prove that if $a \neq 1$ then $b \neq 3$. What are the
      values of $a$ and $b$ if $a \neq 1$? (Remember that $a$ and $b$ are whole
      numbers.)
\end{exlist}
\end{multicols}

\end{document}
